% Provenance (never printed): see sources/outline_sources.md and sources/book1/BRIEFS.md (g9/04).
\chapter{Algorithms on Numbers}\label{ch:g9:algorithms-on-numbers}

Is 1,000,003 a prime number? To find out by hand, you would try to divide it
by 2, by 3, by 4, and so on: about a million divisions, weeks of work. A short
\omterm{def:g4:programs-and-operating-system:program}{program} does them in a moment. A better \omterm{def:g4:programs-and-operating-system:program}{program} does a thousand times fewer,
because it knows where to stop. This chapter \omterm{def:g4:programs-and-operating-system:program}{programs} some of the oldest
\omterm{def:g4:algorithms-repeat-and-decide:algorithm}{algorithms} there are, and counts the work they do.

\begin{recall}
A \omterm{def:g7:functions:function}{function} takes \omterm{def:g7:functions:parameter}{parameters} and returns a \omterm{def:g6:variables-and-input:value}{value}
(\cref{ch:g7:functions}); a \code{while} \omterm{def:g4:algorithms-repeat-and-decide:loop}{loop} repeats as long as a \omterm{def:g4:algorithms-repeat-and-decide:condition}{condition}
holds (\cref{ch:g6:while-loops}); a \omterm{def:g8:searching-and-sorting:binary-search}{binary search} halves the part of a \omterm{def:g8:searching-and-sorting:sorted}{sorted
list} that is left at each step (\cref{ch:g8:searching-and-sorting}). From
mathematics: a divisor of $n$ divides it with no remainder; a prime number has
exactly two divisors, 1 and itself; the greatest common divisor (gcd) of two
numbers is found by Euclid's \omterm{def:g4:algorithms-repeat-and-decide:algorithm}{algorithm}.
\end{recall}

\section{Divisors and primality by trial division}

The operator \texttt{\%} gives the remainder of a division: \texttt{n \% d == 0}
means that \code{d} divides \code{n}. Listing the divisors is one \omterm{def:g4:algorithms-repeat-and-decide:loop}{loop}:

\omcode{code/grade-9/04-algorithms-on-numbers/listings/divisors.py}{1}{11}{The
divisors of a number.}\label{lst:g9:algorithms-on-numbers:divisors}

\Needspace*{6\baselineskip}
\omoutput{code/grade-9/04-algorithms-on-numbers/outputs/divisors.txt}{36 has nine
divisors, 37 only two.}\label{out:g9:algorithms-on-numbers:divisors}

\begin{definition}[Primality test]\label{def:g9:algorithms-on-numbers:primality-test}
A \emph{primality test}\index{primality test} is an \omterm{def:g4:algorithms-repeat-and-decide:algorithm}{algorithm} that tells, for a
whole number, whether it is prime.
\end{definition}

\begin{method}[Trial division]\label{met:g9:algorithms-on-numbers:trial-division}
To test\index{trial division} whether $n \ge 2$ is prime, try to divide it by 2,
3, 4, and so on. If one of them divides $n$, it is not prime; if none does, it
is.
\end{method}

The \omterm{def:g4:programs-and-operating-system:program}{program} below counts the divisions it tries:

\omcode{code/grade-9/04-algorithms-on-numbers/listings/is_prime_slow.py}{1}{12}{\omterm{met:g9:algorithms-on-numbers:trial-division}{Trial
division} by every number below n.}\label{lst:g9:algorithms-on-numbers:slow}

\Needspace*{7\baselineskip}
\omoutput{code/grade-9/04-algorithms-on-numbers/outputs/is_prime_slow.txt}{A prime
costs a division for every number below it.}\label{out:g9:algorithms-on-numbers:slow}

$91 = 7 \times 13$ is found out after 6 divisions. A prime is the worst case:
nothing divides it, so every number from 2 to $n - 1$ is tried, more than a
million divisions for 1,000,003.

\section{Stopping at the square root}

Divisors come in pairs: if $d$ divides $n$, so does $n / d$. For 36, the pairs
are $1 \times 36$, $2 \times 18$, $3 \times 12$, $4 \times 9$ and $6 \times 6$.

\begin{omfigure}
\begin{tikzpicture}[x=0.17cm, y=0.17cm]
  \foreach \w/\h/\y in {36/1/0, 18/2/-3, 12/3/-7, 9/4/-12, 6/6/-18} {
    \fill[ocBlue!25] (0,\y) rectangle ++(\w,-\h);
    \draw[ocInk] (0,\y) rectangle ++(\w,-\h);
    \node[anchor=west, font=\small] at (\w+1,\y-\h/2) {$\h \times \w$};
  }
  \fill[ocRed!30] (0,-18) rectangle ++(6,-6);
  \draw[ocRed, thick] (0,-18) rectangle ++(6,-6);
\end{tikzpicture}

{\small The divisor pairs of 36 as rectangles of area 36. In each pair, one side
is at most 6 and the other at least 6; the square $6 \times 6$ sits in the
middle.}
\end{omfigure}

\begin{proposition}[A divisor below the square root]\label{prop:g9:algorithms-on-numbers:sqrt-bound}
If a whole number $n \ge 2$ is not prime, it has a divisor $d$ with $2 \le d$ and
$d \times d \le n$.
\end{proposition}

\begin{proof}
As $n$ is not prime, $n = d \times e$ with $d$ and $e$ both between 2 and
$n - 1$. Say $d \le e$. If $d \times d$ were larger than $n$, then
$d \times e \ge d \times d > n$, which is false. So $d \times d \le n$.
\end{proof}

So the test can stop as soon as $d \times d > n$: if no divisor was found by
then, there is none.

\omcode{code/grade-9/04-algorithms-on-numbers/listings/is_prime.py}{1}{17}{\omterm{met:g9:algorithms-on-numbers:trial-division}{Trial
division} up to the square root.}\label{lst:g9:algorithms-on-numbers:fast}

\Needspace*{7\baselineskip}
\omoutput{code/grade-9/04-algorithms-on-numbers/outputs/is_prime.txt}{999 divisions
instead of 1,000,001.}\label{out:g9:algorithms-on-numbers:fast}

The square root of 1,000,003 is just above 1,000, so the \omterm{def:g4:algorithms-repeat-and-decide:loop}{loop} tries $d = 2$ to
$1{,}000$: 999 divisions. For a number $n$, the work falls from about $n$
divisions to about $\sqrt{n}$. The test is written \code{d * d <= n} rather than
with a square root, so that it uses whole numbers only.

\section{The sieve for all the primes below a number}

To find \emph{all} the primes below 100, testing each number one by one repeats
work. A \omterm{def:g7:lists:list}{list} can remember what is already known.

\begin{method}[Sieve of Eratosthenes]\label{met:g9:algorithms-on-numbers:sieve}
To find\index{sieve of Eratosthenes} the primes below $n$, write the numbers
from 2 to $n - 1$. Take the smallest number not crossed out: it is prime. Cross
out its multiples. Repeat with the next number not crossed out, until the end of
the \omterm{def:g7:lists:list}{list}.
\end{method}

In Python, a \omterm{def:g7:lists:list}{list} of $n$ Booleans says, for each number, whether it is still in:

\omcode{code/grade-9/04-algorithms-on-numbers/listings/sieve.py}{1}{19}{The
sieve.}\label{lst:g9:algorithms-on-numbers:sieve}

\Needspace*{6\baselineskip}
\omoutput{code/grade-9/04-algorithms-on-numbers/outputs/sieve.txt}{The primes below
50, and how many there are below
100,000.}\label{out:g9:algorithms-on-numbers:sieve}

Crossing out starts at $p \times p$, not at $2 \times p$: Exercise
\ref{exo:g9:algorithms-on-numbers:9} asks why that is enough. Printed as a
\omterm{def:g9:grids-and-images:grid}{grid} of ten rows, the sieve shows where the primes fall:

\omcode{code/grade-9/04-algorithms-on-numbers/listings/sieve_grid.py}{1}{12}{The
primes up to 100, as a \omterm{def:g9:grids-and-images:grid}{grid}.}\label{lst:g9:algorithms-on-numbers:grid}

\omoutput{code/grade-9/04-algorithms-on-numbers/outputs/sieve_grid.txt}{A dot for each
number crossed out.}\label{out:g9:algorithms-on-numbers:grid}

\def\gnineprimes{2,3,5,7,11,13,17,19,23,29,31,37,41,43,47,53,59,61,67,71,73,79,83,89,97}

\begin{omfigure}
\begin{tikzpicture}[x=0.62cm, y=0.62cm]
  \foreach \k in {1,...,100} {
    \pgfmathtruncatemacro{\row}{(\k-1)/10}
    \pgfmathtruncatemacro{\col}{mod(\k-1,10)}
    \draw[black!30, fill=black!4] (\col,-\row) rectangle ++(1,-1);
    \node[font=\scriptsize, black!45] at (\col+0.5,-\row-0.5) {\k};
  }
  \foreach \p in \gnineprimes {
    \pgfmathtruncatemacro{\row}{(\p-1)/10}
    \pgfmathtruncatemacro{\col}{mod(\p-1,10)}
    \draw[ocInk, fill=ocYellow!60] (\col,-\row) rectangle ++(1,-1);
    \node[font=\scriptsize\bfseries] at (\col+0.5,-\row-0.5) {\p};
  }
\end{tikzpicture}

{\small The 25 primes up to 100, from the \omterm{def:g1:what-is-a-computer:input}{output} of the sieve. After 2 and 5,
every prime ends in 1, 3, 7 or 9: the other columns hold multiples of 2 or of 5.}
\end{omfigure}

\begin{history}[A librarian's sieve]
% ledger: g9:eratosthenes
Eratosthenes was born in Cyrene, in today's Libya, about 276~BC, and died in
Alexandria about 194~BC. He was the third librarian of the great library of
Alexandria. He is remembered for the
prime number sieve that carries his name, still used today in a modified form.
\end{history}

How much faster is the sieve? This \omterm{def:g4:programs-and-operating-system:program}{program} finds the primes below 100,000 both
ways and times them:

\omcode{code/grade-9/04-algorithms-on-numbers/listings/primes_timing.py}{1}{15}{\omterm{met:g9:algorithms-on-numbers:trial-division}{Trial
division} against the sieve.}\label{lst:g9:algorithms-on-numbers:timing}

\Needspace*{6\baselineskip}
\omoutput{code/grade-9/04-algorithms-on-numbers/outputs/primes_timing.txt}{Timings on
the author's laptop; yours will differ.}\label{out:g9:algorithms-on-numbers:timing}

Both find the same 9,592 primes; the sieve is many times faster, because each
number is crossed out by a few additions instead of being divided again and
again.

\begin{dated}{2026}
% ledger: g9:largest-prime-2026
The largest prime number known in 2026 is $2^{136\,279\,841} - 1$, found in
October 2024 by a project in which volunteers lend their \omterm{def:g1:what-is-a-computer:computer}{computers} to the
search. Written out, it has 41,024,320 digits. No sieve could reach it: it was
proved prime by a test made for numbers of the form $2^p - 1$.
\end{dated}

\section{Euclid's algorithm programmed}

The gcd of two numbers also divides their difference. Euclid's \omterm{def:g4:algorithms-repeat-and-decide:algorithm}{algorithm} in its
oldest form subtracts the smaller number from the larger until both are equal:

\omcode{code/grade-9/04-algorithms-on-numbers/listings/gcd_subtract.py}{1}{12}{Euclid
by subtraction.}\label{lst:g9:algorithms-on-numbers:subtract}

\Needspace*{11\baselineskip}
\omoutput{code/grade-9/04-algorithms-on-numbers/outputs/gcd_subtract.txt}{The gcd of
84 and 30.}\label{out:g9:algorithms-on-numbers:subtract}

Subtracting 30 from 84 twice is what a division does at once: $84 = 2 \times 30
+ 24$. Euclid's \omterm{def:g4:algorithms-repeat-and-decide:algorithm}{algorithm} by division replaces the pair $(a, b)$ by $(b, a \bmod
b)$, the remainder, until $b$ is 0:

\omcode{code/grade-9/04-algorithms-on-numbers/listings/gcd.py}{1}{12}{Euclid by
division, counting the steps.}\label{lst:g9:algorithms-on-numbers:gcd}

\Needspace*{7\baselineskip}
\omoutput{code/grade-9/04-algorithms-on-numbers/outputs/gcd.txt}{A gcd and the number
of steps.}\label{out:g9:algorithms-on-numbers:gcd}

The line \texttt{a, b = b, a \% b} changes both \omterm{def:g6:variables-and-input:value}{variables} at once: the new
\code{b} is computed from the old \code{a} and \code{b}. Subtraction would need
999,999 steps for the last pair; division needs 2.

\begin{omfigure}
\begin{tikzpicture}[x=0.27cm, y=0.27cm]
  \draw[ocInk, thick] (0,0) rectangle (21,13);
  \foreach \x/\y/\s in {0/0/13, 13/5/8, 16/0/5, 13/0/3, 13/3/2, 15/3/1, 15/4/1} {
    \draw[ocInk, fill=ocBlue!15] (\x,\y) rectangle ++(\s,\s);
  }
  \foreach \x/\y/\s in {0/0/13, 13/5/8, 16/0/5, 13/0/3} {
    \node[font=\small] at (\x+\s/2,\y+\s/2) {\s};
  }
  \node[anchor=west, align=left, font=\small] at (23,6.5) {$21 = 1 \times 13 + 8$\\ $13 = 1 \times 8 + 5$\\
    $8 = 1 \times 5 + 3$\\ $5 = 1 \times 3 + 2$\\ $3 = 1 \times 2 + 1$\\ $2 = 2 \times 1 + 0$};
\end{tikzpicture}

{\small Euclid on 21 and 13: cut the largest square from the rectangle, then from
what is left, again and again. The last squares, 1 by 1, fit exactly: the gcd is
1. Each step cuts as many squares as its quotient: one, five times, then two
squares of side 1; six steps in all.}
\end{omfigure}

\section{Why Euclid is fast}

\begin{proposition}[The remainders halve]\label{prop:g9:algorithms-on-numbers:euclid-halves}
If $0 < b \le a$, then the remainder $a \bmod b$ is less than $a / 2$.
\end{proposition}

\begin{proof}
Two cases. If $b \le a / 2$, the remainder is smaller than $b$, so smaller than
$a / 2$. If $b > a / 2$, then $b$ fits once into $a$ and the remainder is
$a - b$, which is less than $a - a / 2 = a / 2$.
\end{proof}

After two steps, the pair $(a, b)$ has become $(b, r)$, then $(r, \ldots)$:
the larger number is now $r$, less than half of $a$. Every two steps halve the
larger number at least, so for numbers below 1,000, which can be halved only
ten times before reaching 1, Euclid needs at most about twenty steps. The
chart below shows the real counts: they grow very slowly.

\begin{omfigure}
\begin{tikzpicture}
  \begin{axis}[width=12cm, height=5cm, xmin=0, xmax=202, ymin=0, ymax=11,
      xlabel={larger number $a$}, ylabel={most steps},
      ytick={0,2,...,10}, grid=major, grid style={black!10}]
    \addplot[only marks, mark=*, mark size=0.9pt, ocBlue] table[x=a, y=steps]
      {code/grade-9/04-algorithms-on-numbers/figures/out/steps.dat};
    \addplot[only marks, mark=*, mark size=2.2pt, ocRed] coordinates
      {(2,1) (3,2) (5,3) (8,4) (13,5) (21,6) (34,7) (55,8) (89,9) (144,10)};
    \node[anchor=east, font=\small, ocRed, fill=white, inner sep=1pt] at (axis cs:198,1.5) {red: a new record, at $a$ = 2, 3, 5, 8, 13, 21, \ldots};
  \end{axis}
\end{tikzpicture}

{\small For each $a$ up to 200, the largest number of division steps of Euclid's
\omterm{def:g4:algorithms-repeat-and-decide:algorithm}{algorithm} over all pairs $(a, b)$ with $b < a$, computed by the chapter's
\code{gcd}. A new record (red) is set only at a few \omterm{def:g6:variables-and-input:value}{values} of $a$.}
\end{omfigure}

\section{A square root by halving}

The halving of \omterm{def:g8:searching-and-sorting:binary-search}{binary search} also finds numbers. $1 \times 1 < 2 < 2 \times 2$,
so $\sqrt{2}$ lies between 1 and 2.

\begin{method}[Square root by halving]\label{met:g9:algorithms-on-numbers:sqrt-halving}
To approach $\sqrt{2}$, start with an interval whose ends' squares are below
and above 2. Take its middle. If the middle's square is below 2, the middle
becomes the lower end; otherwise it becomes the upper end. Repeat: each step
halves the interval.
\end{method}

\Needspace*{12\baselineskip}
\omcode{code/grade-9/04-algorithms-on-numbers/listings/sqrt_halving.py}{1}{9}{Ten
halvings for the square root of 2.}\label{lst:g9:algorithms-on-numbers:sqrt}

\omoutput{code/grade-9/04-algorithms-on-numbers/outputs/sqrt_halving.txt}{The interval
shrinks around 1.414.}\label{out:g9:algorithms-on-numbers:sqrt}

After ten halvings the interval is $1 / 1024$ wide, less than a thousandth: the
square root of 2 is 1.414 to three decimals. Each halving gains one \omterm{def:g5:bits-and-bytes:binary}{binary}
digit, so about one decimal digit every three or four steps.

\section{Exercises}

\begin{exercise}[$\star$]\label{exo:g9:algorithms-on-numbers:1}
What does \code{divisors(28)} return? Add up the divisors smaller than 28. What
do you notice?
\end{exercise}

\begin{exercise}[$\star$]\label{exo:g9:algorithms-on-numbers:2}
On the numbers from 2 to 30, cross out the multiples of 2, then of 3, as the
sieve does. Which numbers are left? Which of them are not prime?
\end{exercise}

\begin{exercise}[$\star$]\label{exo:g9:algorithms-on-numbers:3}
Write the pairs that \code{gcd(48, 18)} goes through, step by step. How many
steps does it take, and what does it return?
\end{exercise}

\begin{exercise}[$\star$]\label{exo:g9:algorithms-on-numbers:4}
A number is \emph{perfect} when it equals the sum of its divisors smaller than
itself, like 6 $= 1 + 2 + 3$. Write a \omterm{def:g4:programs-and-operating-system:program}{program} that prints the perfect numbers
below 1,000.
\end{exercise}

\begin{exercise}[$\star\star$]\label{exo:g9:algorithms-on-numbers:5}
This test says that 1 is prime. Why? Fix it.

\omcode{code/grade-9/04-algorithms-on-numbers/exercises/is_prime_one_bug.py}{1}{7}{A
\omterm{def:g9:algorithms-on-numbers:primality-test}{primality test} with a gap.}\label{lst:g9:algorithms-on-numbers:one-bug}
\end{exercise}

\begin{exercise}[$\star\star$]\label{exo:g9:algorithms-on-numbers:6}
In this version the \omterm{def:g4:algorithms-repeat-and-decide:loop}{loop} runs while \code{d * d < n}. Which numbers does it
wrongly call prime? Find the smallest one, then fix the test.

\omcode{code/grade-9/04-algorithms-on-numbers/exercises/square_bug.py}{1}{9}{One
\omterm{def:g7:strings:string}{character} changed.}\label{lst:g9:algorithms-on-numbers:square-bug}
\end{exercise}

\begin{exercise}[$\star\star$]\label{exo:g9:algorithms-on-numbers:7}
Write a \omterm{def:g7:functions:function}{function} \code{gcd3(a, b, c)} for the gcd of three numbers, using the
gcd of two. Try it on 84, 126 and 210.
\end{exercise}

\begin{exercise}[$\star\star$]\label{exo:g9:algorithms-on-numbers:8}
Twin primes are two primes that differ by 2, like 11 and 13. Write a \omterm{def:g4:programs-and-operating-system:program}{program}
that counts the pairs of twin primes below 1,000 and prints the last pair.
\end{exercise}

\begin{exercise}[$\star\star$]\label{exo:g9:algorithms-on-numbers:9}
The sieve starts crossing out the multiples of $p$ at $p \times p$. Why have the
smaller multiples, such as $2 \times p$ and $3 \times p$, already been crossed
out?
\end{exercise}

\begin{exercise}[$\star\star\star$]\label{exo:g9:algorithms-on-numbers:10}
10,007 is prime. Without running the \omterm{def:g4:programs-and-operating-system:program}{program}, how many divisions does
\code{is_prime(10007)} try?
\end{exercise}

\begin{exercise}[$\star\star\star$]\label{exo:g9:algorithms-on-numbers:11}
Use the proposition on halving remainders to find an upper bound on the number
of division steps of Euclid's \omterm{def:g4:algorithms-repeat-and-decide:algorithm}{algorithm} for two numbers below 1,000. Compare it
with the real record, found by the weekend problem.
\end{exercise}

\begin{exercise}[$\star\star\star$]\label{exo:g9:algorithms-on-numbers:12}
Write a \omterm{def:g7:functions:function}{function} \code{next_prime(n)} that returns the smallest prime larger
than \code{n}, together with the total number of divisions it tried on the way.
Add a few lines that ask the user for \code{n}.
\end{exercise}

\section{Problem: The Slowest Pairs for Euclid}

\begin{problem}[{Weekend problem --- the slowest pairs for Euclid}]\label{pb:g9:algorithms-on-numbers:1}
Euclid's \omterm{def:g4:algorithms-repeat-and-decide:algorithm}{algorithm} is fast, but some pairs make it work harder than others.
Which pair of numbers below 1,000 takes the most steps?

\medskip\noindent\textbf{Part I --- By hand.}
\begin{enumerate}
  \item Count the division steps of Euclid's \omterm{def:g4:algorithms-repeat-and-decide:algorithm}{algorithm} for $(13, 8)$.
  \item The same for $(21, 13)$; compare with the figure of the squares.
  \item The same for $(100, 1)$. Are large numbers always slow?
\end{enumerate}

\medskip\noindent\textbf{Part II --- A search for records.}

\omcode{code/grade-9/04-algorithms-on-numbers/listings/slowest_pairs.py}{1}{9}{Every
pair below 1,000, printing each new record.}\label{lst:g9:algorithms-on-numbers:slowest}

\begin{enumerate}[resume]
  \item How many pairs does the \omterm{def:g4:programs-and-operating-system:program}{program} try?
  \item Why does it print only some pairs, and not every pair?
  \item Run it. Which pair holds the final record, and with how many steps?
  \item Look at the numbers in the record pairs. How is each number made from
    the two before it?
  \item What is the next record pair, above 1,000?
\end{enumerate}

\medskip\noindent\textbf{Part III --- Why these pairs.}
\begin{enumerate}[resume]
  \item For a record pair $(a, b)$, what is the quotient at each step?
  \item Show that one step of Euclid on a record pair gives the record pair just
    before it.
  \item Why do quotients of 1 make the \omterm{def:g4:algorithms-repeat-and-decide:algorithm}{algorithm} as slow as it can be?
  \item State the result.
\end{enumerate}
\end{problem}
